A Tabela 4.4 fornece as frequências por faixa de 1 dB. Essas frequências são:
\begin{align*}
90\text{–}91 &: 8, & 88\text{–}89 &: 5, & 86\text{–}87 &: 3, & 84\text{–}85 &: 2,\\
82\text{–}83 &: 4, & 80\text{–}81 &: 2, & 78\text{–}79 &: 3, & 76\text{–}77 &: 0,\\
74\text{–}75 &: 6, & 72\text{–}73 &: 4, & 70\text{–}71 &: 6, & 68\text{–}69 &: 3,\\
66\text{–}67 &: 7, & 64\text{–}65 &: 5, & 62\text{–}63 &: 10, & 60\text{–}61 &: 14,\\
58\text{–}59 &: 6, & 56\text{–}57 &: 7, & 54\text{–}55 &: 4, & 52\text{–}53 &: 1, & 50\text{–}51 &: 0.
\end{align*}
Para um histograma de 4 dB de largura, essas observações agrupam-se aproximadamente em:
\begin{align*}
50\text{–}53 &: 1,\\
54\text{–}57 &: 11,\\
58\text{–}61 &: 20,\\
62\text{–}65 &: 15,\\
66\text{–}69 &: 10,\\
70\text{–}73 &: 10,\\
74\text{–}77 &: 6,\\
78\text{–}81 &: 5,\\
82\text{–}85 &: 6,\\
86\text{–}89 &: 8,
\end{align*}
com ainda 8 leituras no intervalo superior 90–91 dB. O histograma deve ser desenhado usando essas frequências como alturas das barras.  